%------------------------------------------------------------------------------%
\begin{question}
    Diagonalize ortogonalmente a matriz
    \[
      B =
      \begin{pmatrix}
        4 & 0 \\
        0 & 4
      \end{pmatrix}
    \]
    Determine as matrizes $P$ e $D$ tais que
    \(
      B = P D P^t
    \)

    \bigskip
    \textit{Essa matriz já é diagonal, este exercício é apenas um treino}
\end{question}

%------------------------------------------------------------------------------%
\begin{resolution}{Autovalores}
  Polinômio característico
  \pause
  \[
    p(\lambda)
    \pause
    = \det(B-\lambda I)
    \pause
    = \begin{vmatrix}
      4-\lambda & 0         \\
      0         & 4-\lambda
    \end{vmatrix}
    \pause
    = (4-\lambda)^{2}
  \]

  \pause
  Autovalor
  \[
    \lambda = 4
  \]
\end{resolution}

%------------------------------------------------------------------------------%
\begin{resolution}{Autoespaço}
\begin{columns}[t]
\column{0.4\textwidth}
  \onslide<+->{Equação característica}
  \begin{align*}
    \onslide<+->{(B-4I)\vec{x}&=\vec{0}}
    \\
    \onslide<+->{\begin{pmatrix}
      0 & 0 \\
      0 & 0
    \end{pmatrix}
    \begin{pmatrix}
        x\\y
    \end{pmatrix}
    &=
    \begin{pmatrix}
        0\\0
    \end{pmatrix}}
  \end{align*}
  \onslide<+->{Não há restrições}
\column{0.4\textwidth}
  \onslide<+->{Autoespaço
  \[
    E(4) = \R^2
  \]}
\end{columns}
\end{resolution}

%------------------------------------------------------------------------------%
\begin{resolution}{Base ortonormal}
  Podemos escolher os vetores ortonormais
  \[
    \pause
    \vec{u}_{1} =
    \begin{pmatrix}
      1 \\
      0
    \end{pmatrix}
    \qquad\qquad
    \pause
    \vec{u}_{2} =
    \begin{pmatrix}
      0 \\
      1
    \end{pmatrix}
  \]
  \pause
  \begin{align*}
    \norm{\vec{u}_{1}}          & =1 \\
    \norm{\vec{u}_{2}}          & =1 \\
    \vec{u}_{1}\cdot\vec{u}_{2} & =0
  \end{align*}
\end{resolution}

%------------------------------------------------------------------------------%
\begin{resolution}{Diagonalização}
  Organizamos os autovetores nas colunas de $P$
  \[
    P =
    \begin{pmatrix}
      1 & 0 \\
      0 & 1
    \end{pmatrix}
  \]

  A matriz diagonal é
  \[
    D =
    \begin{pmatrix}
      4 & 0 \\
      0 & 4
    \end{pmatrix}
  \]

  Portanto
  \[
    B = P D P^t
  \]
\end{resolution}

%------------------------------------------------------------------------------%
